\documentclass[12pt,a4paper]{report}

% ─── Packages ─────────────────────────────────────────────────────────────
\usepackage[utf8]{inputenc}
\usepackage[T1]{fontenc}
\usepackage[french]{babel}
\usepackage{geometry}
\usepackage{graphicx}
\usepackage{hyperref}
\usepackage{listings}
\usepackage{xcolor}
\usepackage{booktabs}
\usepackage{array}
\usepackage{longtable}
\usepackage{amsmath}
\usepackage{float}
\usepackage{caption}
\usepackage{subcaption}
\usepackage{fancyhdr}
\usepackage{titlesec}
\usepackage{enumitem}
\usepackage{mdframed}
\usepackage{tikz}
\usepackage{forest}

% ─── Mise en page ─────────────────────────────────────────────────────────
\geometry{
    top=2.5cm, bottom=2.5cm,
    left=3cm, right=2.5cm
}

% ─── En-têtes et pieds de page ────────────────────────────────────────────
\pagestyle{fancy}
\fancyhf{}
\fancyhead[L]{\small\leftmark}
\fancyhead[R]{\small WIDO — Moteur Sémantique}
\fancyfoot[C]{\thepage}
\renewcommand{\headrulewidth}{0.4pt}

% ─── Couleurs ─────────────────────────────────────────────────────────────
\definecolor{codeblue}{RGB}{0,120,215}
\definecolor{codegray}{RGB}{128,128,128}
\definecolor{codegreen}{RGB}{0,128,0}
\definecolor{codebg}{RGB}{247,247,247}
\definecolor{jdmblue}{RGB}{30,100,200}

% ─── Style code ───────────────────────────────────────────────────────────
\lstdefinestyle{wido}{
    language=JavaScript,
    backgroundcolor=\color{codebg},
    basicstyle=\ttfamily\small,
    keywordstyle=\color{codeblue}\bfseries,
    commentstyle=\color{codegray}\itshape,
    stringstyle=\color{codegreen},
    numberstyle=\tiny\color{codegray},
    numbers=left,
    numbersep=5pt,
    breaklines=true,
    frame=single,
    framesep=4pt,
    rulecolor=\color{codegray},
    tabsize=4,
    showstringspaces=false,
    captionpos=b
}

% Style pour les requêtes WIDO
\lstdefinestyle{widoquery}{
    basicstyle=\ttfamily\normalsize,
    backgroundcolor=\color{codebg},
    frame=single,
    framesep=5pt,
    rulecolor=\color{jdmblue},
    breaklines=true
}

\lstset{style=wido}

% ─── Commandes personnalisées ──────────────────────────────────────────────
\newcommand{\wido}{\textsc{WIDO}}
\newcommand{\jdm}{\textsc{JeuxDeMots}}
\newcommand{\var}[1]{\texttt{\$#1}}
\newcommand{\rel}[1]{\texttt{#1}}
\newcommand{\code}[1]{\texttt{#1}}

% Boîte de mise en valeur
\newenvironment{widobox}{
    \begin{mdframed}[
        linecolor=jdmblue,linewidth=1.5pt,
        backgroundcolor=blue!3,
        skipabove=8pt, skipbelow=8pt,
        innerleftmargin=10pt, innerrightmargin=10pt
    ]
}{
    \end{mdframed}
}

% ─── Hyperref ─────────────────────────────────────────────────────────────
\hypersetup{
    colorlinks=true,
    linkcolor=jdmblue,
    urlcolor=codeblue,
    citecolor=codegray,
    pdftitle={WIDO — Moteur Sémantique pour JeuxDeMots},
    pdfauthor={Mohamed TEDJINI & Maurice DJOBO}
}

% ─── Début du document ────────────────────────────────────────────────────
\begin{document}

% ─── Page de titre ────────────────────────────────────────────────────────
\begin{titlepage}
    \centering
    \vspace*{2cm}
    {\Huge\bfseries WIDO\par}
    {\Large Moteur Sémantique pour JeuxDeMots\par}
    \vspace{0.5cm}
    {\large Langage de Requêtes Formelles avec Heuristiques d'Optimisation\par}
    \vspace{2cm}
    \hrule height 1pt
    \vspace{0.5cm}
    {\large Rapport de projet TER\par}
    {\normalsize Master Informatique — Spécialité ICO\par}
    {\normalsize Université de Montpellier, 2025--2026\par}
    \vspace{0.5cm}
    \hrule height 1pt
    \vspace{2cm}
    {\Large\bfseries Mohamed TEDJINI\par}
    {\Large\bfseries Maurice DJOBO\par}
    \vspace{1.5cm}
    {\normalsize Encadrant : \`a pr\'eciser\par}
    \vfill
    {\normalsize Mai 2026\par}
\end{titlepage}

% ─── Résumé ───────────────────────────────────────────────────────────────
\begin{abstract}
\wido{} est un moteur de recherche sémantique permettant d'interroger le graphe de connaissances \jdm{} (JDM) via un langage de requêtes formel à variables. L'utilisateur formule des requêtes telles que \texttt{(\$x r\_isa animal) ET (\$x r\_has\_part queue)}, que le moteur parse en arbre de décision, optimise par des heuristiques structurelles et de cardinalité, puis exécute en interrogeant l'API JDM v0. Les résultats sont croisés par jointures ensemblistes basées sur les identifiants des nœuds, classés par score de pertinence, et présentés avec leurs preuves dans une interface web glassmorphique.
\end{abstract}

% ─── Table des matières ───────────────────────────────────────────────────
\tableofcontents
\clearpage

% ─── Corps du rapport ─────────────────────────────────────────────────────
\chapter{Introduction}
\label{chap:intro}

Ce projet de Travail d'Étude et de Recherche (TER) s'inscrit dans le cadre du Master Informatique, spécialité ICO. Il porte sur la conception et l'implémentation d'un moteur de recherche sémantique nommé \wido{}, capable d'interroger le graphe de connaissances \jdm{} via un langage de requêtes formel à variables.

\section{Problématique}

Les graphes de connaissances comme \jdm{} contiennent des millions de relations entre concepts. L'enjeu est de permettre à un utilisateur non expert de formuler des requêtes complexes en combinant plusieurs relations, opérateurs logiques et contraintes textuelles, sans avoir à connaître la structure interne de la base.

\section{Objectifs}

Les objectifs de ce projet sont les suivants :
\begin{enumerate}
    \item Définir un langage de requêtes formel adapté à JDM ;
    \item Implémenter un parseur syntaxique produisant un arbre de décision (AST) ;
    \item Concevoir un planificateur d'exécution basé sur des heuristiques structurelles et de cardinalité ;
    \item Réaliser un moteur d'exécution qui interroge l'API JDM v0 et croiser les résultats par jointures sur les identifiants des nœuds ;
    \item Développer une interface web permettant la saisie de requêtes et la visualisation des résultats avec leurs preuves ;
    \item Évaluer les performances du moteur via un benchmark automatique.
\end{enumerate}

\section{Plan du rapport}

Le chapitre~\ref{chap:contexte} présente le contexte du projet et le graphe JeuxDeMots. Le chapitre~\ref{chap:archi} décrit l'architecture globale du moteur. Les chapitres~\ref{chap:parseur} à~\ref{chap:api} détaillent chaque composant. Le chapitre~\ref{chap:eval} présente les résultats du benchmark. Enfin, les chapitres~\ref{chap:limites} et~\ref{chap:conclu} discutent des limites et ouvrent des pistes d'amélioration.

\chapter{Contexte et données}
\label{chap:contexte}

\section{JeuxDeMots}

\jdm{} est un jeu en ligne sérieux développé par le LIRMM (Montpellier) dont le but est de construire collaborativement un réseau lexical du français. Les joueurs associent des mots à des propriétés, des catégories, des synonymes, des actions possibles, etc. Le résultat est un graphe de connaissances contenant des millions de nœuds et de relations.

\section{Structure du graphe JDM}

Chaque \textbf{nœud} représente un concept (mot, expression, syntagme) et est identifié par un entier unique (ID JDM). Chaque \textbf{relation} est un triplet (nœud source, type de relation, nœud cible) associé à un poids (score de confiance) allant de 1 à plusieurs milliers.

Les types de relations principaux utilisés dans \wido{} sont :

\begin{table}[H]
    \centering
    \begin{tabular}{lll}
        \toprule
        Nom & ID & Signification \\
        \midrule
        \rel{r\_isa} & 6 & Appartient à la catégorie \\
        \rel{r\_has\_part} & 9 & Possède la partie \\
        \rel{r\_can\_eat} & 43 & Peut manger \\
        \rel{r\_color} & 42 & A pour couleur \\
        \rel{r\_carac} & 12 & A pour caractéristique \\
        \rel{r\_synt\_synon} & 1 & Synonyme syntaxique \\
        \bottomrule
    \end{tabular}
    \caption{Types de relations JDM utilisés dans WIDO}
    \label{tab:relations}
\end{table}

\section{API JDM v0}

L'API REST publique du LIRMM expose les données JDM via des endpoints structurés. Elle supporte la pagination native (\code{limit} + \code{offset}) et le filtrage par type de relation (\code{types\_ids}) et poids minimum (\code{min\_weight}). Le tableau~\ref{tab:endpoints} résume les endpoints exploités dans \wido{}.

\begin{table}[H]
    \centering
    \begin{tabular}{ll}
        \toprule
        Endpoint & Usage \\
        \midrule
        \code{/v0/node\_by\_name/\{name\}} & Résoudre un nom en nœud JDM \\
        \code{/v0/relations/from/\{name\}} & Récupérer les relations sortantes \\
        \code{/v0/relations/to/\{name\}} & Récupérer les relations entrantes \\
        \code{/v0/relations/from\_by\_id/\{id\}} & Relations sortantes par ID \\
        \code{/v0/relations/to\_by\_id/\{id\}} & Relations entrantes par ID \\
        \code{/v0/relations/from\_by\_id/\{id1\}/to\_by\_id/\{id2\}} & Vérification d'un couple \\
        \bottomrule
    \end{tabular}
    \caption{Endpoints API JDM utilisés par WIDO}
    \label{tab:endpoints}
\end{table}

\chapter{État de l'art}
\label{chap:etat_art}

\section{Moteurs de recherche sémantique}

Les moteurs de recherche sémantique sur des graphes de connaissances ont été largement étudiés dans le contexte du Web sémantique. SPARQL~\cite{sparql2013} est le standard pour interroger des graphes RDF (Resource Description Framework). Des moteurs comme Virtuoso ou GraphDB permettent d'exécuter des requêtes SPARQL sur des bases de connaissances volumineuses.

\section{Spécificité de JDM}

À la différence de DBpedia ou Wikidata, JDM est un graphe informel construit collaborativement. Il n'expose pas nativement de point SPARQL, d'où la nécessité d'une couche d'abstraction comme \wido{}.

\section{Optimisation de requêtes}

L'optimisation de l'ordre d'exécution des jointures est un problème classique en bases de données relationnelles (algorithme de Selinger~\cite{selinger1979}). Dans le contexte des graphes, les heuristiques de cardinalité permettent d'estimer le coût d'une jointure avant de l'exécuter. Notre approche combine une heuristique structurelle (coût syntaxique) et une estimation de cardinalité par sondage de l'API.

\section{Comparaison par identifiant}

Dans un graphe de connaissances, comparer deux nœuds par leur identifiant plutôt que par leur nom est une pratique recommandée pour éviter les ambiguïtés dues aux homonymes et aux variantes typographiques~\cite{noy2009}.

\chapter{Architecture du moteur WIDO}
\label{chap:archi}

\section{Vue d'ensemble}

Le moteur \wido{} suit un pipeline en cinq étapes (figure~\ref{fig:pipeline}) :

\begin{enumerate}
    \item \textbf{Parsing} : la requête textuelle est transformée en arbre de décision (AST) par le parseur ;
    \item \textbf{Estimation de cardinalité} : pour chaque clause, un appel API sonde estime le nombre de résultats ;
    \item \textbf{Planification} : le planificateur trie les clauses ET dans l'ordre d'exécution optimal ;
    \item \textbf{Exécution} : le moteur interroge l'API JDM et croise les résultats par jointures sur les IDs ;
    \item \textbf{Présentation} : les résultats sont triés par score, enrichis de preuves, et affichés.
\end{enumerate}

\begin{figure}[H]
    \centering
    % Schéma ASCII simplifié — remplacer par un diagramme TikZ si disponible
    \begin{verbatim}
        Requête utilisateur
              ↓
         [Parseur] → AST
              ↓
      [Cardinalité] → Map{clause → taille}
              ↓
       [Planificateur] → Plan ordonné
              ↓
     [Moteur d'exécution] ↔ [API JDM + Cache]
              ↓
          Résultats
    \end{verbatim}
    \caption{Pipeline de traitement d'une requête WIDO}
    \label{fig:pipeline}
\end{figure}

\section{Technologies utilisées}

\begin{itemize}
    \item \textbf{Environnement d'exécution} : Node.js v14+
    \item \textbf{Serveur HTTP} : Express.js
    \item \textbf{Client HTTP} : Axios (requêtes vers l'API JDM)
    \item \textbf{Cache disque} : fs-extra, fichiers JSON indexés par hash MD5
    \item \textbf{Interface} : HTML5 + CSS3 (glassmorphique) + JavaScript vanilla
\end{itemize}

\section{Structure des fichiers}

\begin{table}[H]
    \centering
    \begin{tabular}{ll}
        \toprule
        Fichier & Rôle \\
        \midrule
        \code{node.js} & Point d'entrée, routage Express, orchestration \\
        \code{parseur.js} & Analyse syntaxique → AST \\
        \code{heuristiques.js} & Planification par complexité et cardinalité \\
        \code{moteurExecution.js} & Exécution, jointures, intersections \\
        \code{jdmApi.js} & Client API JDM + pagination + cache \\
        \code{cacheManager.js} & Cache disque MD5 persistant \\
        \code{benchmark.js} & Tests automatisés + métriques \\
        \code{verify\_logic.js} & Tests de validation du moteur \\
        \bottomrule
    \end{tabular}
    \caption{Fichiers du projet WIDO}
    \label{tab:fichiers}
\end{table}

\chapter{Parseur syntaxique}
\label{chap:parseur}

\section{Langage de requêtes WIDO}

Le langage \wido{} permet d'exprimer des requêtes formelles à variables sur le graphe JDM. Sa grammaire (simplifiée) est la suivante :

\begin{verbatim}
requête  ::= expr
expr     ::= terme (ET terme)*
terme    ::= facteur (OU facteur)*
facteur  ::= '(' clause ')'
           | '(' expr ')'
clause   ::= '(' sujet relation objet ')'
           | '(' variable '=' filtre ')'
sujet    ::= variable | constante
objet    ::= variable | constante
variable ::= '$' identifiant
filtre   ::= [a-z%]+
\end{verbatim}

\section{Exemples de requêtes}

\begin{lstlisting}[style=widoquery, caption=Exemples de requêtes WIDO]
($x r_isa animal)
($x r_isa animal) ET ($x = ch%)
($x r_isa artiste) ET (($x = ba%) OU ($x = Ba%))
($x r_isa animal) ET ($y r_isa animal) ET ($x r_can_eat $y)
(chat r_has_part $y) ET ($y r_isa $z)
\end{lstlisting}

\section{Structure de l'AST}

Le parseur produit un arbre de décision récursif. Chaque nœud interne représente un opérateur logique (ET/OU), chaque feuille est une clause atomique (relation ou filtre).

\begin{lstlisting}[caption=AST produit pour \texttt{(\$x r\_isa animal) ET (\$x = ch\%)}]
{
  "type": "NOEUD_LOGIQUE",
  "operateur": "ET",
  "gauche": {
    "type": "CLAUSE_RELATION",
    "variable": "$x",
    "relation": "r_isa",
    "cible": "animal"
  },
  "droite": {
    "type": "CLAUSE_FILTRE",
    "variable": "$x",
    "filtre": "ch%"
  }
}
\end{lstlisting}

\section{Gestion des erreurs syntaxiques}

Lorsqu'une clause est incomplète (ex: \texttt{(\$x r\_isa)}), le parseur lève une exception qui est capturée par le serveur. Celui-ci retourne une réponse JSON propre avec \code{statut: "Erreur"} et un message explicatif, sans jamais retourner de HTTP 500 ni de champs \code{undefined}.

\chapter{Heuristiques d'optimisation}
\label{chap:heuristiques}

\section{Principe général}

Pour les requêtes conjonctives (ET), l'ordre d'exécution des clauses influe directement sur le nombre d'appels API et la taille des ensembles intermédiaires. L'objectif du planificateur est de minimiser cet espace de recherche en exécutant d'abord les clauses les plus \textit{sélectives} (qui retournent peu de résultats).

\section{Heuristique structurelle}

La première heuristique est purement syntaxique : elle attribue un score de complexité à chaque clause selon sa structure et les variables déjà ancrées à ce stade de l'exécution.

\begin{table}[H]
    \centering
    \begin{tabular}{clp{7cm}}
        \toprule
        Score & Type de clause & Exemple \\
        \midrule
        0.5 & Filtre sur variable ancrée & \texttt{\$x = ch\%} (après ancrage de \var{x}) \\
        1   & Ancrage par constante & \texttt{(\$x r\_isa animal)} \\
        2   & Vérification (2 vars ancrées) & Jointure finale \\
        5   & Exploration ciblée (1 var ancrée) & \texttt{(\$x r\_can\_eat \$y)} si \var{x} connu \\
        8   & Filtre avant ancrage & \texttt{\$x = ba\%} (avant ancrage) \\
        10  & Jointure sans ancrage & Interdit en première position \\
        \bottomrule
    \end{tabular}
    \caption{Scores de complexité structurelle des clauses WIDO}
    \label{tab:complexity}
\end{table}

\section{Heuristique de cardinalité}

La deuxième heuristique estime la cardinalité (taille de la liste de résultats) de chaque clause d'ancrage par constante en effectuant un appel API sonde avec \code{limit=500}. Si la liste est de taille $n < 500$, la cardinalité est $n$. Sinon, elle est notée $\geq 500$ (grande liste).

\begin{widobox}
\textbf{Exemple d'application :} Pour la requête \texttt{(\$x r\_isa animal) ET (\$x r\_has\_part queue)},\\
si \rel{r\_isa animal} renvoie $\geq 500$ résultats et \rel{r\_has\_part queue} en renvoie 105,\\
le planificateur exécute d'abord \rel{r\_has\_part queue} (plus sélectif), puis filtre les résultats.
\end{widobox}

\section{Mise à jour dynamique des variables ancrées}

Après l'exécution de chaque clause, le planificateur met à jour l'ensemble des variables déjà résolues. Cela modifie le score de complexité des clauses suivantes et permet de détecter qu'une jointure Var→Var peut désormais s'appuyer sur une variable ancrée.

\section{Plan d'exécution (exemple)}

Pour la requête \texttt{(\$x r\_isa animal) ET (\$y r\_isa animal) ET (\$x r\_can\_eat \$y)} :

\begin{table}[H]
    \centering
    \begin{tabular}{clcl}
        \toprule
        Rang & Clause & Complexité & Raison \\
        \midrule
        1 & \texttt{(\$x r\_isa animal)} & 1 & Ancrage initial — cardinalité $\geq 500$ \\
        2 & \texttt{(\$y r\_isa animal)} & 1 & Ancrage initial — cardinalité $\geq 500$ \\
        3 & \texttt{(\$x r\_can\_eat \$y)} & 2 & Vérification ($\var{x}$ et $\var{y}$ ancrés) \\
        \bottomrule
    \end{tabular}
    \caption{Plan d'exécution de la requête officielle à 2 variables}
    \label{tab:plan}
\end{table}

\chapter{Moteur d'exécution}
\label{chap:moteur}

\section{Principe d'exécution}

Le moteur d'exécution parcourt le plan produit par le planificateur et maintient un \textbf{contexte} : une liste de tuples partiels représentant les liaisons variable $\rightarrow$ nœud JDM découvertes jusqu'à présent.

\section{Trois cas de traitement d'une clause de relation}

\subsection{Cas 1 : Ancrage initial (Constante $\rightarrow$ Variable)}

La variable cible n'est pas encore ancrée. Le moteur interroge l'API JDM (\code{/from/} ou \code{/to/}) pour obtenir la liste des nœuds satisfaisant la clause, et crée de nouveaux tuples.

\textit{Exemple :} \texttt{(\$x r\_isa animal)} $\rightarrow$ appel à \code{/v0/relations/to/animal?types\_ids=6}. \\
Résultat : \texttt{[\{\$x: \{id: 39642, name: "chat"\}\}, \{\$x: \{id: ..., name: "chien"\}\}, ...]}

\subsection{Cas 2 : Vérification (Variable ancrée $\rightarrow$ Constante)}

La variable source est déjà ancrée. Le moteur récupère la liste valide depuis l'API, construit un \code{Set} d'IDs, et filtre le contexte en O(1).

\textit{Exemple :} \texttt{(\$x r\_has\_part queue)} après ancrage de \var{x} $\rightarrow$ appel unique + filtrage local.

\subsection{Cas 3 : Jointure (Variable $\rightarrow$ Variable)}

Les deux variables sont des variables, au moins l'une étant déjà ancrée. Le moteur itère sur les candidats de la variable ancrée, effectue un appel API par candidat (via \code{from\_by\_id} ou \code{to\_by\_id}), et conserve les couples valides.

\section{Comparaison par ID JDM}

Toutes les intersections et jointures utilisent les \textbf{identifiants numériques} des nœuds JDM, et non leurs noms textuels. Cette stratégie garantit des comparaisons exactes en O(1) avec un \code{Set<id>}, et évite les ambiguïtés dues aux homonymes ou aux variantes typographiques.

\begin{lstlisting}[caption=Comparaison par ID dans le moteur d'exécution]
// Construction d'un Set d'IDs valides (en O(n))
const validIds = new Set(apiResults.map(r => r.id));

// Filtrage local en O(1) par ID
return contexte.filter(tuple =>
    tuple[$x] && validIds.has(tuple[$x].id)
);
\end{lstlisting}

\section{Scores et preuves}

Chaque tuple porte deux métadonnées :
\begin{itemize}
    \item \code{\_\_score} : somme des poids des relations traversées ;
    \item \code{\_\_preuves} : liste des clauses qui justifient l'inclusion de ce tuple dans le résultat.
\end{itemize}

Les résultats finaux sont triés par score décroissant.

\section{Traitement des opérateurs OU}

Pour les branches OU, le moteur exécute chaque branche indépendamment et calcule l'union des résultats. Les doublons (même combinaison d'IDs) sont éliminés en conservant le tuple de score maximal.

\chapter{Client API et cache}
\label{chap:api}

\section{Pagination de l'API JDM}

L'API JDM retourne au plus $n$ résultats par appel. Pour récupérer des listes complètes, \wido{} implémente une boucle de pagination native avec les paramètres \code{limit} et \code{offset}.

\begin{lstlisting}[caption=Boucle de pagination (extrait simplifié de jdmApi.js)]
let offset = 0, page = 0;
while (true) {
    if (page >= maxPages) break;         // Limite de sécurité
    if (results.length >= maxTotal) break;

    const url = `${base}?types_ids=${relId}
                 &limit=${limit}&offset=${offset}`;
    const data = await apiCall(url);

    results.push(...data.relations);
    page++; offset += limit;

    if (data.relations.length < limit) break; // Derniere page
}
\end{lstlisting}

Les paramètres de sécurité sont : \code{limit=1000}, \code{maxPages=5}, \code{maxTotal=5000}.

\section{Cache disque MD5}

Le cache disque évite les appels redondants à l'API JDM. Chaque clé de cache est un hash MD5 de l'URL ou des paramètres de la requête. Les données sont stockées en JSON sur le disque.

\begin{table}[H]
    \centering
    \begin{tabular}{ll}
        \toprule
        Type de cache & Clé \\
        \midrule
        Résolution de nœud & \code{node\_by\_name:\{name\}} \\
        Relations paginées & \code{relsPaged:\{dir\}:\{node\}:\{rel\}:...} \\
        Relations par ID & \code{relsPaged:\{dir\}:id:\{nodeId\}:\{rel\}:...} \\
        Cardinalité estimée & \code{cardinality:\{dir\}:\{node\}:\{rel\}:...} \\
        Vérification de couple & \code{check:\{id1\}:\{rel\}:\{id2\}} \\
        \bottomrule
    \end{tabular}
    \caption{Structure des clés de cache WIDO}
    \label{tab:cache}
\end{table}

\section{Gestion des erreurs API}

\begin{itemize}
    \item \textbf{HTTP 404} : nœud ou relation introuvable → résultat vide (cas légitime)
    \item \textbf{HTTP 500} : erreur serveur JDM → retry 1× après 500ms
    \item \textbf{Timeout ($>$15s)} : appel abandonné, warning remonté à l'interface
    \item \textbf{Erreur réseau} : log et continuation avec les données disponibles
\end{itemize}

La résilience du moteur repose sur ces mécanismes : une clause qui échoue ne bloque pas les autres.

\chapter{Interface utilisateur}
\label{chap:interface}

\section{Design glassmorphique}

L'interface web adopte un style glassmorphique : fond sombre, éléments translucides avec \code{backdrop-filter: blur()}, typographie claire.

\section{Zone de saisie et résultats}

L'utilisateur saisit sa requête. La touche Entrée déclenche la recherche. Les résultats affichent : variables liées, score de pertinence, et preuves (bouton toggle).

\section{Panneau de diagnostic}

L'interface affiche : statut, warnings, joinStats (candidats testés, couples trouvés), pagination (pages et relations récupérées), AST.

\section{Gestion des erreurs}

Quand \code{statut: "Erreur"}, bannière rouge avec message. Aucun \code{undefined} affiché.

\chapter{Évaluation et benchmark}
\label{chap:eval}

\section{Protocole}

Le benchmark automatique exécute 14 requêtes représentatives sur le serveur en cours d'exécution et mesure : durée, appels API, cache, pagination, joinStats, et résultats.

\section{Requêtes testées}

\begin{longtable}{cll}
    \toprule
    \# & Requête & Description \\
    \midrule
    1 & \texttt{(\$x r\_isa animal)} & Simple, ancrage \\
    2 & \texttt{(\$x r\_isa animal) ET (\$x = ch\%)} & Filtre texte \\
    3 & \texttt{(\$x r\_isa animal) ET (\$x r\_has\_part queue)} & Jointure ET \\
    4 & \texttt{((\$x r\_isa mammifere) OU (\$x r\_isa oiseau)) ET (\$x = ch\%)} & OU imbriqué \\
    5 & \texttt{(\$x r\_isa artiste) ET ((\$x = ba\%) OU (\$x = Ba\%))} & Requête officielle \\
    6 & \texttt{(\$x r\_isa animal) ET ((\$x r\_has\_part aile) OU ...)} & ET + OU \\
    7 & \texttt{(\$x r\_isa animal) ET (\$y r\_isa animal) ET (\$x r\_can\_eat \$y)} & 2 variables \\
    8 & \texttt{(chat r\_isa \$x)} & Direction inverse \\
    9 & \texttt{(chat r\_has\_part \$x)} & Exploration \\
    10 & \texttt{(chat r\_has\_part \$y) ET (\$y r\_isa \$z)} & 3 variables \\
    11 & \texttt{(\$x r\_isa animal) ET (\$x r\_has\_part \$y) ET (\$y = pa\%)} & 3 vars + filtre \\
    12 & \texttt{(lion r\_can\_eat \$y) ET (\$y r\_isa animal)} & Ciblée \\
    13 & \texttt{(\$x r\_isa)} & Erreur syntaxique \\
    14 & \texttt{(\$x relation\_inconnue animal)} & Relation invalide \\
    \bottomrule
\end{longtable}

\section{Analyse des résultats}

Les résultats du benchmark sont générés automatiquement dans \code{benchmark\_results.md} après exécution de \code{node benchmark.js}. Voir ce fichier pour les métriques exactes (durées, appels API, joinStats).

Les observations générales sont :
\begin{itemize}
    \item Les requêtes simples (1-2 variables) sont traitées en moins de 2s ;
    \item Le cache réduit significativement les durées des requêtes répétées ;
    \item La requête officielle à 2 variables (r\_can\_eat) peut retourner 0 résultat si la relation est peu renseignée dans JDM ;
    \item La pagination est utilisée pour les grandes listes (ex: r\_isa animal $\geq$ 500 entrées).
\end{itemize}

\chapter{Limites et perspectives}
\label{chap:limites}

\section{Limites actuelles}

\paragraph{Densité de JDM.} La qualité des résultats dépend directement de la densité du graphe JDM pour les relations utilisées. La relation \rel{r\_can\_eat} est peu renseignée, ce qui explique que la requête officielle à 2 variables peut retourner peu ou aucun résultat.

\paragraph{Combinatoire des jointures.} Pour les requêtes à 2-3 variables, le nombre de candidats à tester peut exploser. Les limites de sécurité (\code{joinCandidateLimit=500}, \code{joinEarlyStop=1000}) contrôlent cet espace mais peuvent tronquer les résultats.

\paragraph{Polysémie non gérée.} Un mot comme "chat" peut correspondre à plusieurs nœuds JDM (animal, messagerie...). Le moteur prend le nœud canonique sans désambiguïsation.

\paragraph{Pagination bornée.} La pagination est limitée à 5000 relations par clause pour protéger l'API du LIRMM. Des relations plus volumineuses sont tronquées.

\section{Perspectives d'amélioration}

\begin{itemize}
    \item \textbf{Désambiguïsation} : exploiter l'endpoint \code{/refinements/} pour gérer la polysémie ;
    \item \textbf{Index local} : construire un index partiel des relations fréquentes pour réduire les appels API ;
    \item \textbf{Jointures parallèles} : tester plusieurs candidats en parallèle avec \code{Promise.all} contrôlé ;
    \item \textbf{Requêtes récursives} : explorer des chemins de longueur variable dans le graphe ;
    \item \textbf{Langage naturel} : ajouter une couche de traduction langage naturel → requête WIDO.
\end{itemize}

\chapter{Conclusion}
\label{chap:conclu}

Ce projet a conduit à la conception et à l'implémentation complète du moteur sémantique \wido{}, capable d'interroger le graphe de connaissances \jdm{} via un langage de requêtes formel à variables.

Les contributions principales sont :
\begin{itemize}
    \item Un \textbf{parseur} robuste gérant les opérateurs ET, OU, les filtres textuels et les erreurs syntaxiques ;
    \item Un \textbf{planificateur} combinant heuristique structurelle et estimation de cardinalité pour optimiser l'ordre d'exécution ;
    \item Un \textbf{moteur d'exécution} utilisant les IDs JDM pour des jointures exactes et efficaces ;
    \item Une \textbf{pagination API} permettant de récupérer des listes complètes sans surcharger le service ;
    \item Une \textbf{interface glassmorphique} affichant proprement les résultats, les scores, les preuves et les statistiques de diagnostic ;
    \item Un \textbf{benchmark automatique} de 14 requêtes générant des métriques reproductibles.
\end{itemize}

Le moteur est pleinement fonctionnel et a été validé sur les requêtes officielles du sujet TER. Ses limites principales (densité de JDM, combinatoire des jointures) sont inhérentes à la nature du graphe source et non à l'architecture du moteur.

\wido{} constitue une base solide pour des travaux futurs sur la recherche sémantique en langage naturel dans des graphes de connaissances collaboratifs.


% ─── Bibliographie ────────────────────────────────────────────────────────
\bibliographystyle{plain}
\bibliography{bibliographie}

\end{document}
